\paragraph{Exponential observables \(\operatorname{Tr}\mathrm e^{\mathrm ik^\mu A_\mu}\) at \(N=2\) (bounded obstruction, partly numerical).}
Keep Definition~(A), the contours \(\mathcal R_u\) and the displayed limit identity of the paragraph on the regulated Lorentzian integral at finite \(N\). Let \(O_k(A)=\operatorname{Tr}\mathrm e^{\mathrm ik^\mu A_\mu}\) with \(k\in\mathbb C^{10}\), at \(\gamma=0\), regulator (A), contours \(\mathcal R_u\) and \(\epsilon\to0^+\) last. For each fixed \(\epsilon>0\) and \(u<2/3\) the integral converges absolutely and is entire in \(k\), because on \(\mathcal R_u\) the growth of \(O_k\) is at most \(\mathrm e^{c(u,k)\|\tilde A\|}\) against a Gaussian regulator, and the fixed-\(\epsilon\) Stokes identity extends to it. At \(N=2\) the matrix \(k^\mu A_\mu\) has eigenvalues \(\pm\rho_X/2\) with \(\rho_X^2=\sum_{a=1}^{3}(k^\mu A^a_\mu)^2\), so \(O_k=2\cos(\rho_X/2)\) is branch-independent, and \(|2\cos(\omega/2)|^2=4(\cos^2(\operatorname{Re}\omega/2)+\sinh^2(\operatorname{Im}\omega/2))\). Along rank-one rays \(\tilde A_\mu=\rho_\mu H\), \(\operatorname{Tr}H^2=1\), the Euclidean weight \(\mathrm e^{-S_{\mathrm E}}\) is not exponentially damped and tends to a bounded factor along the corner (\(S_{\mathrm E}\) vanishes on the exact rank-one rays), the Pfaffian vanishes on the ray and is polynomial in the coordinates nearby, and the remaining factors are polynomial (numerical, in a reduced chart), while \(|O_k|\ge 2\sinh\!\bigl(b_u(\hat\rho;k)\rho/2\bigr)\) with \(b_u(\hat\rho;k)=|\operatorname{Im}(k^0\mathrm e^{-3\pi\mathrm iu/8}\hat\rho_0+k^i\mathrm e^{\pi\mathrm iu/8}\hat\rho_i)|\). At every \(u\in(0,1]\) the functional \(\hat\rho\mapsto\operatorname{Im}(k^0\mathrm e^{-3\pi\mathrm iu/8}\hat\rho_0+k^i\mathrm e^{\pi\mathrm iu/8}\hat\rho_i)\) is nonzero for every nonzero real \(k\) (at \(u=1\) its coefficients are \(-\sin(3\pi/8)k^0\) and \(\sin(\pi/8)\vec k\)), so the Euclidean side of the displayed limit identity diverges in modulus for every nonzero real \(k\); more generally it diverges for every \(k\) outside the slice on which \(k^0\mathrm e^{-3\pi\mathrm i/8}\) and all \(k^i\mathrm e^{\pi\mathrm i/8}\) are real, where \(|O_k|\le2\) and that side converges by the bound along the flat directions with polynomial degree \(0\). For \(\operatorname{Im}k\ne0\) such that some commuting unit \(\hat\rho\) has \(|\operatorname{Im}(k\cdot\hat\rho)|>|\operatorname{Re}(k\cdot\hat\rho)|\), the commuting-ray contribution to the \(u=0\) integral has \(S_{\mathrm b}=0\) and grows along a sequence \(\epsilon\to0^+\) like \(\mathrm e^{c_2(k)/\epsilon}\) with \(c_2(k)>0\) (commuting-ray model; no transverse stationary-phase bound is claimed), and for real \(k\ne0\) the integrand at \(u=0\) carries \(|O_k|\le2\) pointwise, but no fixed-\(u\) domination in \(\epsilon\) is available at any \(u\in(0,2/3)\): the fixed-\(\epsilon\) modulus on \(\mathcal R_u\) grows without bound as \(\epsilon\to0\) (numerical, for a large-radius \(\cosh\) proxy of \(|O_k|\) in a reduced chart, on the sampled range \(\epsilon\in[0.01,1]\); the rank-one corner estimate of the paragraph on the slice bridge band, applied along a ray with \(b_u(\hat\rho;k)\ne0\), gives growth at least like \(\mathrm e^{c/\epsilon}\) up to polynomial factors). Hence within the class stated the displayed limit identity fails for every nonzero real \(k\) and for every non-real \(k\) outside the slice on which \(k^0\mathrm e^{-3\pi\mathrm i/8}\) and all \(k^i\mathrm e^{\pi\mathrm i/8}\) are real, holds at \(k=0\), and is unresolved on that slice at \(k\ne0\): the displayed identity has no analyticity strip in \(k\) and no non-polynomial extension at \(\gamma=0\) off that slice, within the hypothesis class \(N=2\), regulator (A), contours \(\mathcal R_u\), \(\epsilon\to0^+\) last. This concerns the right-hand side of the displayed identity and the modulus route, not the existence of the limit: by (v) of the paragraph on linear contours that mix the indices, at \(N=2\) the limit of \(\langle O_k\rangle_\epsilon\) exists and equals \(F_{\mathrm E}(\mathrm e^{\mathrm i\pi/4}k^\mu k_\mu)\) for every real spacelike \(k\) and for every \(k\) with \(\Lambda^{\mathrm T}k\) real on \(\mathcal R_{u'}\), \(0<u'<2/3\), although for some of these \(k\) the commuting-ray contribution grows. The obstruction is the massless flat direction, and at fixed \(\gamma<0\) it is absent: the mass contributes on \(\mathcal R_u\) the real exponent \((\gamma/2g_{\mathrm{YM}}^2)[\sin(3\pi u/4)\operatorname{Tr}\tilde A_0^2+\sin(\pi u/4)\sum_{i=1}^{9}\operatorname{Tr}\tilde A_i^2]\), a Gaussian damping on every component that dominates \(\mathrm e^{c(k)\|\tilde A\|}\) uniformly for \(\epsilon\le\epsilon_0<|\gamma|/(2g_{\mathrm{YM}}^2)\), and the displayed limit identity then extends to \(O_k\) for every \(k\in\mathbb C^{10}\) (derived in the paragraph on the \(\gamma\to0^-\) link below). At fixed \(\gamma>0\) the \(u=1\) integral diverges in modulus even at \(k=0\) (numerical), so the contour family and the thimbles of~\cite{ChouNishimuraTripathi2025} act on different contour integrals and produce no competing value; the classicalization of~\cite{AsanoNishimuraPiensukYamamori2025} is tied there to a Lorentz-breaking infrared cutoff and bears neither on this paragraph nor on the statements above. Everything in this paragraph is for the hypothesis class stated in it, does not address \(N\to\infty\) taken first, and derives no signature, real-time response, background or continuation; Assumption~\ref{ass:lorentzian} is not discharged.

\emph{The slice at \(k\ne0\): refinement of the window and model-level growth control.} On the slice \(k^0\mathrm e^{-3\pi\mathrm i/8}=r^0\in\mathbb R\), \(k^i\mathrm e^{\pi\mathrm i/8}=r^i\in\mathbb R\) the identity \(\rho_X^2=\sum_{a=1}^{3}{(k^\mu A^a_\mu)}^2=-4\det(k^\mu A_\mu)\) is exact at \(N=2\), and on \(\mathcal R_1\) the contraction \(k^\mu A_\mu\) reduces to the Hermitian real-linear form \(r^\mu\tilde A_\mu\), so \(\rho_X\) is real and \(|O_k|\le2\) pointwise there; right-hand-side boundedness on the slice is structural. The commuting-ray alternative also refines inside the slice: at \(u=0\), writing the growth functional as \(B=\mathbf v\mathbf v^{\mathrm T}-\mathbf u\mathbf u^{\mathrm T}\) with \(\mathbf u=(\cos(3\pi/8)r^0,\cos(\pi/8)\vec r\,)\), \(\mathbf v=(\sin(3\pi/8)r^0,-\sin(\pi/8)\vec r\,)\), the matrix \(B\) is supported on \(\operatorname{span}\{\mathbf e_0,\hat{\mathbf r}\}\), has rank at most two, and \(\lambda_{\max}(B)=[{(r^0)}^2-r^2+\sqrt{{(r^0)}^4+6{(r^0)}^2r^2+r^4}]/(2\sqrt2)\), which is positive iff \(r^0\ne0\); for pure spatial slice momenta (\(r^0=0\)) there is no commuting ray with growth, the neutral set is \(\hat{\mathbf r}\cdot\hat\rho=0\), and \(B\) is negative off it. For \(r^0\ne0\) the one-dimensional branch integral \(\int_{\mathbb R}\mathrm e^{-Qa^2}\,2\cos(\gamma_{\mathrm c}|a|)\,\mathrm da=2\sqrt{\pi/Q}\,\mathrm e^{-\gamma_{\mathrm c}^2/(4Q)}\), \(\gamma_{\mathrm c}=\mathrm e^{3\pi\mathrm i/8}r^0/2\) frozen at its value on the Hermitian contour, gives in that frozen-coefficient model (with \(g_{\mathrm{YM}}=1\)) pointwise modulus growth \(\mathrm e^{{(r^0)}^2\cos(3\pi u/4-\pi/4)/(16\epsilon)}\) for the soft temporal mode at \(\gamma=0\), saturating as \(\epsilon\to0^+\) at fixed \(\gamma<0\) to a constant of order \(\mathrm e^{{(r^0)}^2\sin(3\pi u/4-\pi/4)/(8|\gamma|)}\), which diverges as \(\gamma\to0^-\) when \(3\pi u/4-\pi/4>0\) and decays when it is negative. In the one-dimensional branch model the coefficient of \(\tilde A_0\) in \(k^\mu A_\mu\) on \(\mathcal R_u\) at pure temporal slice momentum is \(\mathrm e^{3\pi\mathrm i(1-u)/8}r^0\), not its value at \(u=0\); with this contour-dependent coefficient in \(\gamma_{\mathrm c}\) the soft-mode growth exponent at \(\gamma=0\) is \({(r^0)}^2/(16\sqrt2\,\epsilon)\) and the fixed-\(\gamma\) saturation exponent is \(-g_{\mathrm{YM}}^2{(r^0)}^2/(8\sqrt2|\gamma|)\), both independent of \(u\), so the saturation factor above and its growth as \(\gamma\to0^-\) for \(3\pi u/4-\pi/4>0\) are properties of the frozen-coefficient model, not of the contour family; with the contour-dependent coefficient the exponential factor of the signed one-dimensional integral tends at fixed \(\gamma<0\) to a constant below one for every \(u\in(0,2/3)\). In the reduced \(N=2\) integral (regulator (A), contours \(\mathcal R_u\), \(\epsilon\to0^+\) last) with the factorized surrogate of this insertion along the soft temporal mode, coefficient frozen as above (it coincides with the contour-dependent one at \(u=0\)), at pure temporal slice momenta (the exact insertion, evaluated at \(u=0\) for \(\epsilon\ge0.1\) only, gives ratios between \(2.0\) and \(4.1\)), the modulus ratio \(|Z_\epsilon[O_k]|/|Z_\epsilon[1]|\) stays of order one and tends to \(2\) along the ladder down to \(\epsilon=3\cdot10^{-3}\) (fit slope \(0.0000\) of \(\ln\)-ratio against \(1/\epsilon\), against the pointwise corner value \(0.044/\epsilon\) at \(r^0=1\)), and at \(u=0.6\) with \(\gamma=-0.2,-0.05\) it stays between \(0.74\) and \(1.78\): in the reduction neither the commuting-ray growth \(\mathrm e^{c_2(k)/\epsilon}\) nor the pointwise saturation factor of the frozen-coefficient model materialises, consistent with phase-space suppression of the corner branch growth near the soft corner by Vandermonde and measure powers, though no domination lemma is claimed. These are reduced-model computations with oscillation-limited resolution at \(\gamma=0\); whether the left-hand side is controlled on the slice at \(k\ne0\) with \(r^0\ne0\) remains open at full \(N=2\), the \(\gamma\to0^-\) bridge from fixed-\(\gamma<0\) success remains open, and the pure-spatial half of the slice has no commuting ray with growth.
